\chapter{Trading à haute fréquence sur les plateformes centralisées de crypto-actifs}\label{ch:hf:crypto-high-frequency-trading-on-centralised-venues}

Le bitcoin se négocie sur des dizaines de bourses, vingt-quatre heures sur vingt-quatre, sans bande consolidée. Un teneur de marché cote sur l'une, couvre sa position sur une autre, verse le funding
toutes les huit heures sur une troisième, et doit tenir le tout dans le budget de requêtes de chaque plateforme. Dans la semaine simulée de ce chapitre, un teneur de marché qui cote un
contrat perpétuel à deux points de base autour du cours spot d'une autre plateforme encaisse 575\,000~\$ de fourchette, paie 464\,000~\$ pour couvrir sa position et perçoit 144\,000~\$ de
rabais de maker~; une cascade de liquidations ajoute 20\,800~\$ si elle absorbe les ventes forcées à la bonne profondeur. Un budget de requêtes qui n'autorise qu'une recotation toutes les
quatre secondes lui coûte 102\,000~\$ sur cette semaine.

\section{Market making inter-bourses et prise de liquidité}

Les plateformes crypto (livre~3, chapitres~15 et~16) sont des carnets d'ordres distincts, avec une garde des actifs, des frais et des règles propres~; la même crypto-monnaie se négocie sur toutes
elles, et rien n'achemine un ordre de l'une vers le meilleur prix sur l'autre. Makarov et Schoar ont mis en évidence de grandes opportunités d'arbitrage
récurrentes entre plateformes, nettement plus fortes d'un pays à l'autre que dans un même pays, et en ont attribué la persistance aux contrôles des capitaux et à la difficulté de transférer les
fonds~; au sein d'un même pays, les prix restent proches parce que des traders comme ceux que décrit ce chapitre les y maintiennent.

\begin{definition}[Market making inter-bourses]\label{def:hf:crypto-high-frequency-trading-on-centralised-venues:xmm}
Le \emph{market making inter-bourses}\index{market making inter-bourses} consiste à coter sur une plateforme à des prix dérivés du carnet d'une autre plateforme, et
en couvrant chaque exécution aussitôt par une prise de liquidité sur l'autre plateforme, de sorte que le teneur de marché encaisse l'écart entre la fourchette qu'il cote et le coût de
la couverture.
\end{definition}

\begin{definition}[Plateforme de couverture]\label{def:hf:crypto-high-frequency-trading-on-centralised-venues:hedge}
Une \emph{plateforme de couverture}\index{plateforme de couverture} est la bourse sur laquelle un teneur de marché inter-bourses compense les positions que ses cotations accumulent sur
une autre~: en général le carnet le plus profond pour la crypto-monnaie, celui où ses prises de liquidité déplacent le moins le cours et dont le cours lui sert de fair value.
\end{definition}

\begin{omfigure}
\begin{tikzpicture}[font=\footnotesize,>=stealth,box/.style={draw,rounded corners,fill=omThm!8,minimum width=2.6cm,minimum height=0.8cm,align=center}]
\node[box] (a) at (0,0) {plateforme de cotation\\contrat perpétuel};
\node[box,fill=omDef!10] (m) at (5.3,0) {teneur de marché\\(fair = spot)};
\node[box] (b) at (10.6,0) {plateforme de couverture\\carnet spot};
\node[box,fill=gray!10] (c) at (0,-2.1) {clients, arbitragistes,\\ventes forcées};
\node[box,fill=gray!10] (fu) at (5.3,-2.1) {funding\\toutes les 8 heures};
\draw[<->] (m) -- node[above,font=\scriptsize] {cotations, 2 pb} (a);
\draw[->] (b) -- node[above,font=\scriptsize] {prix juste} (m);
\draw[->] (m.south east) to[bend right=15] node[below right] {couverture~: prise} (b.south west);
\draw[->] (c) -- node[left] {exécutions} (a);
\draw[->] (fu) -- (m);
\node[font=\scriptsize,align=center] at (5.3,1.2) {budget de requêtes~: une recotation\\toutes les $1/r$ secondes};
\end{tikzpicture}

\omcaption{Un teneur de marché inter-bourses~: il cote le contrat perpétuel sur une plateforme autour du \omterm{def:hf:fair-value:fair}{prix juste} lu dans le carnet spot d'une autre, et il couvre
chaque exécution là-bas, en prenant de la liquidité, reçoit ou verse le funding sur sa position perpétuelle toutes les huit heures, et ne peut modifier ses cotations qu'aussi souvent que son
budget de requêtes le lui permet. Source~: ce chapitre.}\label{fig:hf:crypto-high-frequency-trading-on-centralised-venues:flow}
\end{omfigure}

La semaine du chapitre (la \cref{fig:hf:crypto-high-frequency-trading-on-centralised-venues:flow})~: une crypto-monnaie à 60\,000~\$ avec une volatilité quotidienne de 3~\%~; les clients
envoient des ordres de 10\,000~\$ sur le perpétuel, espacés d'une demi-seconde en moyenne, dont 52~\% sont des achats~; les arbitragistes prennent 20\,000~\$ chaque fois qu'une cotation est devenue périmée
de plus que sa demi-fourchette, augmentée d'un point de base. Le teneur de marché cote à 2 points de base de part et d'autre du \omterm{def:hf:fair-value:fair}{prix juste}, paie des frais de maker de $-0.5$
point de base (un rabais), et couvre sur le spot en prenant de la liquidité à une demi-fourchette de 0.5 point de base, plus des frais de 1 point de base et 20 millisecondes de slippage. Il détient
au plus 200\,000~\$ d'inventaire perpétuel de chaque côté (la couverture le rend neutre en delta~; la limite borne le risque de base).

\section{Funding et base à grande vitesse}

Un contrat perpétuel (livre~3, chapitre~17) n'a pas d'échéance~; un paiement de funding à dates fixes le maintient proche de l'indice. Un teneur de marché dont
la position perpétuelle, couverte sur le spot, est exposé non pas à la crypto-monnaie mais à la base entre les deux et au funding versé sur la jambe perpétuelle. Dans le
modèle, la légère préférence des clients pour l'achat laisse le teneur de marché en position vendeuse sur le perpétuel (acheteur sur le spot) la plupart de la semaine~; avec un funding de 1 point de base
par intervalle en moyenne, il perçoit quelques centaines de dollars. Le funding compte davantage pour le choix du moment du teneur de marché que pour son niveau~: les positions qui franchissent
un horodatage de funding paient ou encaissent le taux de l'intervalle, celles ouvertes après ne sont pas concernées, et ce taux est connu à l'avance grâce à l'indice de prime.

\section{Flux de liquidation}

\begin{definition}[Absorption des liquidations]\label{def:hf:crypto-high-frequency-trading-on-centralised-venues:absorb}
L'\emph{absorption des liquidations}\index{absorption des liquidations} consiste à fournir de la liquidité aux ordres forcés que le moteur de liquidation d'une plateforme envoie dans son
carnet pendant une cascade, à des prix qui escomptent la vente forcée, avec une couverture sur une autre plateforme et la capacité de tenir la position jusqu'à ce que la
dislocation se résorbe.
\end{definition}

La cascade injectée~: le quatrième jour à midi, le perpétuel chute de 80 points de base sous le spot en cinq minutes, puis se redresse pendant l'heure suivante, avec une constante de temps de dix minutes
(la \cref{fig:hf:crypto-high-frequency-trading-on-centralised-venues:cascade})~; des ventes forcées de 50\,000~\$ par seconde arrivent pendant cinq minutes. Un teneur de marché qui
retire son cours acheteur pendant que les ventes forcées arrivent ne gagne rien sur elles (600~\$ sur le compte de la cascade de la semaine). Celui qui les absorbe avec une
limite de 5 millions de dollars dès la première seconde achète tôt, à des décotes faibles, et gagne 5\,200~\$ sur la remontée de la base~; celui qui attend que le perpétuel soit
à 40 points de base sous le spot gagne 20\,800~\$~; à 60 points de base, il est trop tard pour remplir la limite et gagne 17\,800~\$.

\begin{omfigure}
\begin{tikzpicture}
\begin{axis}[width=11cm,height=5.4cm,axis y line*=right,axis x line=none,xmin=-2,xmax=60,ymin=-90,ymax=10,
  ylabel={\small perpétuel moins spot (pb)},tick label style={font=\footnotesize},table/col sep=comma]
\addplot[gray,thick] table[x=minute,y=basis]{figdata/hft/24-crypto-high-frequency-trading-on-centralised-venues/cascade.csv};
\label{plot:hf:cbasis}
\end{axis}
\begin{axis}[width=11cm,height=5.4cm,axis y line*=left,xmin=-2,xmax=60,xlabel={\small minutes depuis le début de la cascade},
  ylabel={\small position perpétuelle (millions de dollars)},ymin=-0.5,ymax=5.5,tick label style={font=\footnotesize},grid=major,grid style={gray!25},
  table/col sep=comma,legend style={font=\scriptsize,at={(0.5,-0.3)},anchor=north},legend columns=3]
\addplot[omThm,thick] table[x=minute,y=absorb]{figdata/hft/24-crypto-high-frequency-trading-on-centralised-venues/cascade.csv};
\addlegendentry{absorption à partir de $-40$ pb}
\addplot[omDef,thick,dashed] table[x=minute,y=pull]{figdata/hft/24-crypto-high-frequency-trading-on-centralised-venues/cascade.csv};
\addlegendentry{retrait du cours acheteur}
\addlegendimage{/pgfplots/refstyle=plot:hf:cbasis}\addlegendentry{base (à droite)}
\end{axis}
\end{tikzpicture}

\omcaption{La cascade de liquidation injectée~: la base du perpétuel par rapport au spot (axe de droite) et la position perpétuelle du teneur de marché, couverte sur le spot (axe de
gauche), pour un teneur de marché qui absorbe les ventes forcées dès que la base atteint 40 points de base, à hauteur de 5 millions de dollars, et pour celui qui retire son cours acheteur. Données~:
\texttt{hf\_crypto.absorption}.}\label{fig:hf:crypto-high-frequency-trading-on-centralised-venues:cascade}
\end{omfigure}

\begin{dated}{2026-09}{Limites, horloges de funding et cascade d'octobre 2025}\label{dat:hf:crypto-high-frequency-trading-on-centralised-venues:limits}
Le point d'accès aux informations sur l'échange d'une grande plateforme spot a renvoyé, le 24 septembre 2026, des limites de poids de requête de 6\,000 par minute et des limites d'ordres de 100 par
10 secondes et 200\,000 par jour. La même plateforme règle le funding de ses contrats à terme à 00:00, 08:00 et 16:00 UTC, avec une composante d'intérêt de 0.01~\% par
intervalle pour la plupart des contrats. Le 10 octobre 2025, après une annonce tarifaire, 19~milliards de dollars de positions crypto à effet de levier ont été liquidés en 24
heures, selon le décompte d'un rapport de recherche d'un éditeur d'actualités crypto (livre~3, chapitre~18).
\end{dated}

\section{Cotation tenant compte des limites de débit}

Une plateforme crypto mesure les requêtes (livre~3, chapitre~15)~: poids de requête par minute, ordres par dix secondes, ordres par jour. La règle la plus contraignante fixe
à quelle fréquence le teneur de marché peut modifier ses cotations. Avec les limites de l'encadré daté et un seul instrument, le plafond quotidien d'ordres autorise 2.3 ordres par seconde, soit 1.16
recotations de deux ordres chacune~: c'est le plafond quotidien, et non celui de dix secondes, qui est contraignant. Coter cinq instruments dans le même budget laisse à chacun 0.23 recotation par
seconde.

\begin{omfigure}
\begin{tikzpicture}
\begin{axis}[width=11cm,height=5.2cm,xlabel={\small recotations par seconde},ylabel={\small milliers de \$ sur la semaine},xmode=log,log basis x=10,
  xtick={0.1,0.2,0.25,0.5,1},xticklabels={0.1,0.2,0.25,0.5,1},xmin=0.09,xmax=1.1,ymin=-700,ymax=320,tick label style={font=\footnotesize},
  grid=major,grid style={gray!25},table/col sep=comma,legend style={font=\scriptsize,at={(0.97,0.05)},anchor=south east}]
\addplot[omThm,thick,mark=*,mark size=1.4pt] table[x=rate,y=total]{figdata/hft/24-crypto-high-frequency-trading-on-centralised-venues/budget.csv};
\addplot[omDef,thick,dashed,mark=square*,mark size=1.3pt] table[x=rate,y=adverse]{figdata/hft/24-crypto-high-frequency-trading-on-centralised-venues/budget.csv};
\draw[gray,dashed] (axis cs:0.09,0) -- (axis cs:1.1,0);
\legend{total de la semaine,prélevé par les arbitragistes}
\end{axis}
\end{tikzpicture}

\omcaption{Le total du teneur de marché sur la semaine simulée et la part perdue au profit des arbitragistes qui saisissent les cotations périmées, en fonction de la fréquence à laquelle son budget de requêtes
lui permet de modifier ses cotations (échelle logarithmique)~: une fois par seconde, toutes les deux, quatre, cinq et dix secondes. Données~:
\texttt{hf\_crypto.budget}.}\label{fig:hf:crypto-high-frequency-trading-on-centralised-venues:budget}
\end{omfigure}

Le coût d'un budget serré tient aux cotations périmées (\cref{fig:hf:crypto-high-frequency-trading-on-centralised-venues:budget})~: à une recotation par seconde, les
arbitragistes ne trouvent rien~; toutes les deux secondes, ils prélèvent 6\,400~\$ sur la semaine~; toutes les quatre, 135\,000~\$~; toutes les dix, 675\,000~\$, et la semaine perd
266\,000~\$. Coter cinq instruments dans le budget d'un seul compte fait passer la perte hebdomadaire de 275\,000~\$ à environ 100\,000~\$~; un teneur de marché ajoute des comptes, les paliers
de teneur de marché des plateformes, assortis de limites plus élevées, ou de moins d'instruments, avant d'ajouter des cotations.

\section{Programmes de teneur de marché et leurs obligations}

Les plateformes rémunèrent les teneurs de marché pour leur liquidité au moyen de programmes (livre~3, chapitre~25)~: une remise sur le volume du teneur de marché si ses cotations respectent une exigence de
disponibilité, soit la part des instants échantillonnés où les deux côtés sont à l'intérieur d'une fourchette maximale et affichent une taille minimale. Le teneur de marché du chapitre cote des deux côtés
et reste à moins de 10 points de base 90~\% du temps~; ses limites d'inventaire, qui retirent un côté lorsqu'il est plein, sont ce qui lui fait perdre sa disponibilité, et
l'exigence du programme est une contrainte pour sa gestion des risques autant que pour sa technologie. Les émetteurs de jetons signent des accords de market making assortis d'obligations de
même nature (livre~3, chapitre~24).

\section{Fiches de stratégie}

\begin{strategyfile}{Market making inter-bourses}\label{strat:hf:crypto-high-frequency-trading-on-centralised-venues:xmm}
\sfield{Qui paie le gain, et pourquoi} Des traders sur la plateforme la moins liquide, qui y paient la fourchette pour une immédiateté que la plateforme profonde fournit à moindre coût.
\sfield{Instruments et plateformes} Le contrat perpétuel ou le spot d'un même actif numérique sur une plateforme de cotation~; le carnet spot le plus profond comme \omterm{def:hf:crypto-high-frequency-trading-on-centralised-venues:hedge}{plateforme de couverture}.
\sfield{Signal} Le carnet de la \omterm{def:hf:crypto-high-frequency-trading-on-centralised-venues:hedge}{plateforme de couverture} comme fair value~; la file d'attente et le flux de la plateforme de cotation.
\sfield{Dimensionnement et exécution} Coter à une demi-fourchette au-dessus du coût de la couverture~; couvrir chaque exécution sans délai~; plafonner l'inventaire pour borner le risque de base.
\sfield{Coûts} La couverture (1.5 points de base par exécution dans le modèle), les transferts entre plateformes, la garde des actifs.
\sfield{Comment elle s'éteint} Plateformes qu'on ne peut pas atteindre avec des fonds~: les primes transfrontalières de Makarov et Schoar ont persisté parce que le capital ne pouvait pas circuler.
\sfield{Horizon, capacité, infrastructure} Des millisecondes~; des comptes et un inventaire sur les deux plateformes.
\sfield{Backtest honnête} Les deux carnets sur une même horloge~; les délais de transfert~; le prix d'exécution réel de la couverture.
\sfield{Sources} Makarov et Schoar (2020)~; ce chapitre.
\end{strategyfile}

\begin{strategyfile}{Prise de liquidité inter-bourses par latence}\label{strat:hf:crypto-high-frequency-trading-on-centralised-venues:taking}
\sfield{Qui paie le gain, et pourquoi} Des teneurs de marché sur des plateformes plus lentes dont les cotations retardent sur le carnet de référence (chapitre~8).
\sfield{Instruments et plateformes} Un même actif numérique sur plusieurs plateformes~; la plateforme leader, en général la plus profonde.
\sfield{Signal} Le mouvement du leader face aux cotations d'une plateforme en retard, au-delà des frais.
\sfield{Dimensionnement et exécution} Prendre la cotation périmée~; compenser sur le leader ou tenir jusqu'à la convergence dans les limites d'inventaire.
\sfield{Coûts} Frais du preneur de liquidité sur les deux plateformes~; le risque que le retard soit une information et non une cotation périmée.
\sfield{Comment elle s'éteint} Des teneurs de marché qui recotent plus vite~: dans le modèle, un teneur de marché qui recote chaque seconde ne laisse rien~; toutes les quatre secondes, 135\,000~\$ par semaine.
\sfield{Horizon, capacité, infrastructure} Des millisecondes~; colocation là où les plateformes la proposent.
\sfield{Backtest honnête} Les cotations de la plateforme en retard telles qu'elles étaient, y compris le calendrier de recotation propre au teneur de marché.
\sfield{Sources} Chapitre~8~; ce chapitre.
\end{strategyfile}

\begin{strategyfile}{Trading sur les horodatages de funding}\label{strat:hf:crypto-high-frequency-trading-on-centralised-venues:funding}
\sfield{Qui paie le gain, et pourquoi} Des traders à effet de levier qui détiennent des contrats perpétuels à travers les moments de funding et paient le taux.
\sfield{Instruments et plateformes} Contrats perpétuels et leurs couvertures en spot~; funding à heures fixes (00:00, 08:00 et 16:00 UTC sur une grande plateforme).
\sfield{Signal} Le taux de financement prévu à partir de l'indice de prime avant chaque horodatage.
\sfield{Dimensionnement et exécution} Tenir le côté qui reçoit à travers l'horodatage, couvert en spot~; débouclér la position après.
\sfield{Coûts} Les frais des deux jambes, payés deux fois~; la base bouge autour de l'horodatage à mesure que d'autres font de même.
\sfield{Comment elle s'éteint} Encombrement (crowding)~: la base bouge avant l'horodatage et emporte le taux avec elle.
\sfield{Horizon, capacité, infrastructure} Quelques minutes autour de chaque horodatage.
\sfield{Backtest honnête} Taux de financement tels que publiés avant l'horodatage~; la base dans les minutes qui l'entourent.
\sfield{Sources} Le livre~3, chapitre~17.
\end{strategyfile}

\begin{strategyfile}{Absorption des flux de liquidation}\label{strat:hf:crypto-high-frequency-trading-on-centralised-venues:liquidation}
\sfield{Qui paie le gain, et pourquoi} Des traders à effet de levier dont le moteur de liquidation vend les positions quel que soit le prix.
\sfield{Instruments et plateformes} Contrats perpétuels pendant les cascades~; spot sur une autre plateforme comme couverture.
\sfield{Signal} La profondeur de la base, les ordres forcés du moteur de liquidation, les positions ouvertes près des prix de liquidation.
\sfield{Dimensionnement et exécution} Absorber une fois la base profonde (40 points de base dans le modèle), jusqu'à la capacité~; couvrir en spot~; tenir jusqu'à la reprise.
\sfield{Coûts} La couverture sur un marché en baisse~; le risque que la dislocation soit le nouveau prix~; la solvabilité propre de la plateforme.
\sfield{Comment elle s'éteint} Des cascades plus profondes que la capacité, et des plateformes qui appliquent l'auto-deleveraging (ADL) aux positions gagnantes des absorbeurs, comme certaines l'ont fait le 10 octobre 2025
lorsque 19~milliards de dollars de positions ont été liquidées en une journée.
\sfield{Horizon, capacité, infrastructure} De quelques minutes à quelques heures~; bilan sur plusieurs plateformes.
\sfield{Backtest honnête} Flux de liquidation tels que le moteur les a envoyés~; règles d'auto-deleveraging en vigueur à la date.
\sfield{Sources} Le livre~3, chapitre~18~; l'encadré daté.
\end{strategyfile}

\begin{strategyfile}{Cotation tenant compte des limites de débit dans le cadre d'un programme}\label{strat:hf:crypto-high-frequency-trading-on-centralised-venues:ratelimit}
\sfield{Qui paie le gain, et pourquoi} La remise du programme de la plateforme, et les clients par la fourchette.
\sfield{Instruments et plateformes} Plateformes de cotation avec budgets de requêtes et programmes de teneur de marché.
\sfield{Signal} La valeur d'une recotation~: le mouvement du \omterm{def:hf:fair-value:fair}{prix juste} depuis la dernière cotation, comparé à la demi-fourchette.
\sfield{Dimensionnement et exécution} Dépenser les requêtes là où la cotation est la plus périmée~; garder les deux côtés affichés pour la disponibilité~; moins d'instruments par compte.
\sfield{Coûts} Cotations périmées lorsque le budget est contraignant (102\,000~\$ par semaine avec une recotation toutes les quatre secondes dans le modèle).
\sfield{Comment elle s'éteint} Budgets réduits par la plateforme, programmes modifiés.
\sfield{Horizon, capacité, infrastructure} Des secondes~; un limiteur de requêtes (\texttt{firm.ratelimit}).
\sfield{Backtest honnête} Les limites de la plateforme à la date (encadré daté), et non les valeurs d'exemple de la documentation.
\sfield{Sources} Le livre~3, chapitres~15 et~25.
\end{strategyfile}

\begin{strategyfile}{Market making sur la base spot-perpétuel}\label{strat:hf:crypto-high-frequency-trading-on-centralised-venues:basis}
\sfield{Qui paie le gain, et pourquoi} Des traders qui paient le levier du perpétuel, par le funding et la base.
\sfield{Instruments et plateformes} Un contrat perpétuel et son spot, sur une ou deux plateformes.
\sfield{Signal} La base comparée à son niveau attendu, compte tenu du taux de financement.
\sfield{Dimensionnement et exécution} Coter la base comme un seul instrument~; couvrir les jambes~; garder la position en carry à travers les moments de funding lorsque le taux rémunère.
\sfield{Coûts} Les frais des deux jambes~; dislocations de la base dans les cascades.
\sfield{Comment elle s'éteint} Retournements du funding et cascades~; le premier perpétuel, le XBTUSD de BitMEX, coté en 2016 avec un funding toutes les huit heures, a fixé le modèle
que toutes les plateformes ont copié (livre~3, chapitre~17).
\sfield{Horizon, capacité, infrastructure} Des heures à des semaines.
\sfield{Backtest honnête} Historiques de funding et de base par plateforme.
\sfield{Sources} Le livre~3, chapitre~17.
\end{strategyfile}

\section{Tutoriel~: 24 heures sur 24}\label{tut:hf:crypto-high-frequency-trading-on-centralised-venues:tut}

\begin{tutorial}
\textbf{Objectif.} Faire tourner une semaine de \omterm{def:hf:crypto-high-frequency-trading-on-centralised-venues:xmm}{market making inter-bourses} avec du funding, une cascade et un budget de requêtes, puis décomposer le résultat.
\textbf{État final~:} les trois figures et le tableau ci-dessous.
\begin{tutsteps}
\item \textbf{Les cotations} ne bougent que lorsque le budget le permet~; les limites d'inventaire, et l'absorption dans une cascade.
\omcode{code/firm/cryptohft/firm_cryptohft.py}{87}{100}{Recotation selon le calendrier du budget~; limite de chaque côté~; absorption seulement lorsque la base est profonde.}\label{lst:hf:crypto-high-frequency-trading-on-centralised-venues:quote}
\item \textbf{Les exécutions et la couverture}~: clients, ventes forcées, arbitragistes sur les cotations périmées~; chaque exécution couverte en spot sans délai.
\omcode{code/firm/cryptohft/firm_cryptohft.py}{101}{121}{Exécutions contre les cotations, edge du teneur de marché par rapport au mid du perpétuel, frais et coût de la couverture.}\label{lst:hf:crypto-high-frequency-trading-on-centralised-venues:fills}
\item \textbf{La semaine}, le balayage des budgets et les variantes d'absorption (\texttt{hf\_crypto})~; le budget issu du limiteur de \texttt{firm.ratelimit}.
\end{tutsteps}
\textbf{Pour aller plus loin.} Introduire un skew dans les cotations selon l'inventaire au lieu de retirer un côté~; laisser le \omterm{def:hf:fair-value:fair}{prix juste} provenir de deux \omterm{def:hf:crypto-high-frequency-trading-on-centralised-venues:hedge}{plateformes de couverture}~; ajouter des délais
de transfert et un gel des retraits sur la plateforme de cotation.
\end{tutorial}

\begin{center}
\small
\begin{tabular}{@{}lr@{}}
\toprule
composante de la semaine (recotation chaque seconde, absorption à partir de 40 pb) & \$\\
\midrule
fourchette encaissée sur le perpétuel & 574\,912\\
perte due aux cotations périmées & 0\\
couverture sur le spot (demi-fourchette, frais, glissement) & $-464\,310$\\
rabais aux teneurs de marché & 143\,728\\
financement & 274\\
cascade (rattrapage de la base) & 20\,842\\
total & 275\,445\\
\bottomrule
\end{tabular}
\end{center}

\section{Construction~: le teneur de marché crypto inter-bourses}\label{bld:hf:crypto-high-frequency-trading-on-centralised-venues:cryptohft}

\begin{build}
\bfield{Rôle} Simuler le \omterm{def:hf:crypto-high-frequency-trading-on-centralised-venues:xmm}{market making inter-bourses} avec couverture, financement, une cascade de liquidations et un budget de requêtes.
\bfield{Interface} \texttt{Week(seed, days, sigma\_day, clients\_per\_s, clip, buy\_share, cascade\_day, cascade\_depth\_bp, forced\_per\_s)},
\texttt{run(week, half\_bp, requote\_per\_s, inv\_limit, absorb, absorb\_limit, absorb\_from\_bp, record)} et paramètres de frais et de latence,
\texttt{budget\_from(governor, instruments)}. Construit à partir de \texttt{firm.perp}, \texttt{firm.ratelimit} et \texttt{firm.mmprogram} du livre~3.
\bfield{Règles} Dollars~; points de base~; pas d'une seconde~; chaque exécution couverte aussitôt~; financement versé par le détenteur à chaque échéance de huit heures.
\bfield{Tests d'acceptation} \texttt{code/firm/cryptohft/tests/}~: les parties s'additionnent au total~; les cotations fraîches ne sont jamais frappées par les arbitragistes~; les cotations périmées coûtent
davantage et utilisent moins de requêtes~; le budget issu du plafond quotidien du limiteur~; l'absorption rapporte davantage sur la cascade que le retrait et constitue une
position de plusieurs millions.
\bfield{Pour aller plus loin} Skew selon l'inventaire~; deux \omterm{def:hf:crypto-high-frequency-trading-on-centralised-venues:hedge}{plateformes de couverture}~; délais de transfert~; auto-deleveraging de l'absorbeur.
\end{build}

\begin{omsources}
\item I.~Makarov, A.~Schoar, Trading and arbitrage in cryptocurrency markets, \emph{Journal of Financial Economics} 135(2), 2020, 293--319.
\item Le livre~3, chapitres 15, 17, 18 et 25, et leurs sources pour l'encadré daté.
\end{omsources}

\section{Exercices}

\begin{exercise}[$\star$]\label{exo:hf:crypto-high-frequency-trading-on-centralised-venues:1}
Un teneur de marché cote 2 points de base de part et d'autre, reçoit un rabais de teneur de marché de 0.5 point de base, et se couvre en franchissant
une demi-fourchette de 0.5 point de base et supporte des frais de 1 point de base. Quel est son edge par aller-retour, avant glissement~?
\end{exercise}

\begin{exercise}[$\star$]\label{exo:hf:crypto-high-frequency-trading-on-centralised-venues:2}
Une plateforme autorise 200\,000 ordres par jour. Combien de recotations (annulation et nouvel ordre) par seconde cela représente-t-il pour un instrument, et pour cinq~?
\end{exercise}

\begin{exercise}[$\star$]\label{exo:hf:crypto-high-frequency-trading-on-centralised-venues:3}
Un teneur de marché est vendeur de 150\,000~\$ de contrat perpétuel à un horodatage de financement, au taux de 0.01~\%. Que reçoit-il ou que paie-t-il~?
\end{exercise}

\begin{exercise}[$\star\star$]\label{exo:hf:crypto-high-frequency-trading-on-centralised-venues:4}
Pourquoi l'absorption dès la première seconde de la cascade rapporte-t-elle moins qu'une attente jusqu'à ce que la base atteigne 40 points de base de profondeur~?
\end{exercise}

\begin{exercise}[$\star\star$]\label{exo:hf:crypto-high-frequency-trading-on-centralised-venues:5}
Pourquoi un budget de requêtes serré fait-il perdre de l'argent au profit des arbitragistes alors que la fourchette du teneur de marché reste inchangée~?
\end{exercise}

\begin{exercise}[$\star\star$]\label{exo:hf:crypto-high-frequency-trading-on-centralised-venues:6}
Qu'est-ce qui a fait persister les primes entre pays de Makarov et Schoar~?
\end{exercise}

\begin{exercise}[$\star\star\star$]\label{exo:hf:crypto-high-frequency-trading-on-centralised-venues:7}
\emph{Programmation.} Relancer l'absorption avec un seuil de 60 points de base et avec une limite de 2~millions de dollars dès la première seconde. Laquelle des deux fait mieux, et pourquoi~?
\end{exercise}

\begin{exercise}[$\star\star\star$]\label{exo:hf:crypto-high-frequency-trading-on-centralised-venues:8}
\emph{Trouver la faille.} «~Notre backtest cote 20 instruments sur un seul compte et recote chacun chaque seconde~; il a rapporté 5~millions de dollars le mois dernier.~»
\end{exercise}

\section{Problème~: autour de la montre}

\begin{problem}[{Problème du week-end --- autour de la montre}]\label{pb:hf:crypto-high-frequency-trading-on-centralised-venues:1}
Un teneur de marché crypto cote un contrat perpétuel sur une plateforme et se couvre sur une autre pendant une semaine qui contient une cascade de liquidations.

\medskip\noindent\textbf{Partie I --- La structure.}
\begin{enumerate}
\item Définir le \omterm{def:hf:crypto-high-frequency-trading-on-centralised-venues:xmm}{market making inter-bourses} et la \omterm{def:hf:crypto-high-frequency-trading-on-centralised-venues:hedge}{plateforme de couverture}.
\item Qu'ont établi Makarov et Schoar~?
\item Résumer l'encadré daté.
\item Décrire la semaine du chapitre.
\end{enumerate}

\medskip\noindent\textbf{Partie II --- Financement et cascades.}
\begin{enumerate}[resume]
\item À quoi le teneur de marché reste-t-il exposé une fois couvert~?
\item Définir l'\omterm{def:hf:crypto-high-frequency-trading-on-centralised-venues:absorb}{absorption des liquidations}.
\item Donner le résultat de la cascade pour un retrait, une absorption immédiate, et une absorption à partir de 40 points de base ou de 60 points de base.
\item Que peut-il mal tourner avec l'absorption~?
\end{enumerate}

\medskip\noindent\textbf{Partie III --- Budgets et programmes.}
\begin{enumerate}[resume]
\item Laquelle des limites de l'encadré daté est contraignante, et à combien de recotations par seconde~?
\item Donner le résultat de la semaine avec une recotation toutes les 1, 2, 4 et 10 secondes.
\item Que mesure l'exigence de disponibilité d'un programme, et qu'est-ce qui fait perdre de la disponibilité au teneur de marché ici~?
\item Comment répartir un budget serré~?
\end{enumerate}

\medskip\noindent\textbf{Partie IV --- Le verdict.}
\begin{enumerate}[resume]
\item Énoncer le \emph{résultat nommé}~: le P\&L du teneur de marché inter-bourses par composante sur une semaine avec une cascade, et la capture à laquelle il renonce pour rester
dans son budget de requêtes.
\item Quelle composante est la plus importante, et laquelle est la plus fragile~?
\item Qu'apporterait une deuxième \omterm{def:hf:crypto-high-frequency-trading-on-centralised-venues:hedge}{plateforme de couverture}~?
\item Quelle fiche de stratégie est la plus exposée à la défaillance d'une plateforme~?
\item Comment les horodatages de financement doivent-ils modifier la politique d'inventaire du teneur de marché~?
\item Que doit le teneur de marché à la plateforme dans le cadre d'un programme~?
\item Comment tester la cascade du modèle par rapport au record d'octobre 2025~?
\item En une phrase~: que vend un teneur de marché inter-bourses~?
\end{enumerate}
\end{problem}

\section{Questions d'entretien}

\begin{interviewq}[$\star$ \iqroles{trader}]\label{iq:hf:crypto-high-frequency-trading-on-centralised-venues:1}
Le bitcoin vaut 30~dollars de plus sur une plateforme que sur une autre. Que vérifiez-vous avant de le trader~?
\end{interviewq}

\begin{interviewq}[$\star\star$ \iqroles{développeur}]\label{iq:hf:crypto-high-frequency-trading-on-centralised-venues:2}
Quel limiteur de requêtes concevriez-vous pour trois plateformes aux limites différentes, avec un moteur de cotation qui ne doit jamais être banni~?
\end{interviewq}

\begin{interviewq}[$\star\star$ \iqroles{chercheur}]\label{iq:hf:crypto-high-frequency-trading-on-centralised-venues:3}
Comment prévoiriez-vous le prochain taux de financement, et quelle part en est connue avant l'horodatage~?
\end{interviewq}

\begin{interviewq}[$\star\star$ \iqroles{risque}]\label{iq:hf:crypto-high-frequency-trading-on-centralised-venues:4}
Votre société détient 20~millions de dollars de cryptomonnaies sur une plateforme qui suspend les retraits. Qu'est-ce qui est exposé, et quelles auraient dû être les limites~?
\end{interviewq}

\begin{interviewq}[$\star\star$ \iqroles{trader}]\label{iq:hf:crypto-high-frequency-trading-on-centralised-venues:5}
Les positions ouvertes sont élevées et la base baisse vite. Placeriez-vous un cours acheteur sur le perpétuel~? Pour quel montant, et où couvririez-vous~?
\end{interviewq}

\begin{interviewq}[$\star\star\star$ \iqroles{chercheur}]\label{iq:hf:crypto-high-frequency-trading-on-centralised-venues:6}
Avec une recotation toutes les $\tau$ secondes, une volatilité des prix $\sigma$ par racine de seconde et une demi-fourchette $h$, estimeriez-vous le taux
auquel un arbitragiste trouve une cotation périmée de plus de $h$.
\end{interviewq}
